% Experiment Designs poster — NEMI 2026, v2
% Six experiment designs grounded in the mechanistic validity paper's
% actual criteria (C6, I6, I9, I10, I11/I12, E5).
% All source disciplines match the paper.
%
%   cd poster && lualatex designs_poster_v2.tex

\documentclass[25pt, a0paper, portrait]{tikzposter}

\usepackage[T1]{fontenc}
\usepackage{amsmath,amssymb}
\usepackage{array}
\usepackage{booktabs}
\usepackage{enumitem}
\usepackage{microtype}

\tikzposterlatexaffectionproofoff

% ── Color palette ───────────────────────────────────────
\definecolor{darktext}{HTML}{2C3E50}
\definecolor{blockbg}{HTML}{FFFFFF}
\definecolor{hookred}{HTML}{C0392B}
\definecolor{designblue}{HTML}{2980B9}
\definecolor{constructgreen}{HTML}{27AE60}
\definecolor{internalblue}{HTML}{2980B9}
\definecolor{externalpurple}{HTML}{8E44AD}

% ── Theme ───────────────────────────────────────────────
\usetheme{Default}

\definecolorstyle{DesignColors}{
  \definecolor{colorOne}{HTML}{2C3E50}
  \definecolor{colorTwo}{HTML}{3498DB}
  \definecolor{colorThree}{HTML}{ECF0F1}
}{
  \colorlet{backgroundcolor}{colorThree}
  \colorlet{framecolor}{colorOne}
  \colorlet{titlefgcolor}{white}
  \colorlet{titlebgcolor}{colorOne}
  \colorlet{blocktitlebgcolor}{designblue}
  \colorlet{blocktitlefgcolor}{white}
  \colorlet{blockbodybgcolor}{blockbg}
  \colorlet{blockbodyfgcolor}{darktext}
  \colorlet{notefgcolor}{darktext}
  \colorlet{notebgcolor}{designblue!12}
  \colorlet{notefrcolor}{designblue}
}
\usecolorstyle{DesignColors}

\defineblockstyle{DesignBlock}{
  titlewidthscale=1.0, bodywidthscale=1.0, titleleft,
  titleoffsetx=0pt, titleoffsety=0pt,
  bodyoffsetx=0pt, bodyoffsety=0pt, bodyverticalshift=0pt,
  roundedcorners=12, linewidth=2pt, titleinnersep=10pt, bodyinnersep=14pt
}{
  \draw[rounded corners=\blockroundedcorners, color=blocktitlebgcolor,
        fill=blockbodybgcolor, line width=\blocklinewidth]
    (blockbody.south west) rectangle (blocktitle.north east);
  \ifBlockHasTitle
    \draw[rounded corners=\blockroundedcorners,
          fill=blocktitlebgcolor, color=blocktitlebgcolor,
          line width=\blocklinewidth]
      (blocktitle.south west) rectangle (blocktitle.north east);
  \fi
}
\useblockstyle{DesignBlock}

\definetitlestyle{DesignTitle}{
  width=\paperwidth, roundedcorners=0, linewidth=0pt, innersep=20pt,
  titletotopverticalspace=15mm, titletoblockverticalspace=25mm
}{
  \draw[fill=titlebgcolor, color=titlebgcolor]
    (\titleposleft, \titleposbottom)
    rectangle (\titleposright, \titlepostop);
}
\usetitlestyle{DesignTitle}

\title{%
  \parbox{0.92\linewidth}{\centering\bfseries
    Experiment Designs We Haven't Run Yet}}
\author{Elliot Tower\quad\texttt{elliot@elliottower.ai}}
\institute{}

\begin{document}

\maketitle[width=\paperwidth]

% ════════════════════════════════════════════════════════
% ROW 1 — THE PITCH + COMPLEMENTATION (C6)
% ════════════════════════════════════════════════════════
\begin{columns}
\column{0.45}

{
\colorlet{blocktitlebgcolor}{hookred}
\block{Six Criteria, Four Source Disciplines}{
  Cardiac stenting for stable angina was standard for three
  decades---500,000 procedures a year in the US. The first trial
  comparing stenting against a matched control (ORBITA, 2017)
  found no significant improvement in exercise tolerance. The
  intervention was real. The causal attribution to the specific
  intervention had never been tested against a procedure that
  matched everything except the active ingredient.

  \vspace{5mm}
  Circuit ablation studies face the same gap. We audit 16
  mechanistic claims against 36 validity criteria drawn from
  philosophy of science, psychometrics, genetics, neuroscience,
  medical microbiology, and pharmacology. Six of these criteria
  map to concrete experiment designs that are rarely or never
  used in mechanistic interpretability.

  \vspace{5mm}
  \begin{center}
  \renewcommand{\arraystretch}{1.3}
  {\small
  \begin{tabular}{@{}llll@{}}
    \toprule
    \textbf{ID} & \textbf{Design} & \textbf{Source} & \textbf{Cost} \\
    \midrule
    C6  & Complementation  & Genetics         & 2 passes \\
    I6  & Double dissoc.   & Neuroscience     & 1 cond. \\
    I9  & Epistatic inter. & Genetics         & 1 cond. \\
    I10 & Rescue           & Genetics         & 1 pass \\
    I11/12 & Onset/offset  & Microbiology     & Varies \\
    E5  & Dose--response   & Pharmacology     & Free \\
    \bottomrule
  \end{tabular}}
  \end{center}
}
}

\column{0.55}

{
\colorlet{blocktitlebgcolor}{constructgreen}
\block{C6: Complementation --- Same Unit or Different?}{
  In genetics, two mutations that both break a function might be in
  the same gene or in different genes. Benzer's cis-trans test (1955)
  distinguishes them: place one mutation on each chromosome (the
  \emph{trans} configuration). If function recovers, the mutations
  are in different genes---they \emph{complement} each other.

  \vspace{5mm}
  \textbf{The MI translation.} S-inhibition heads and negative name
  movers both suppress the repeated name in IOI. To test whether
  they are the same functional unit: run two forward passes, one
  with each group intact. Construct a hybrid activation by taking
  the S-inhibition contribution from pass~1 and the negative
  name-mover contribution from pass~2 (the trans configuration).
  If the model recovers, they serve independent roles. If the
  hybrid fails like either single ablation, they are the same
  functional unit.

  \vspace{5mm}
  \begin{center}
  \renewcommand{\arraystretch}{1.3}
  {\large
  \begin{tabular}{@{}lll@{}}
    \toprule
    \textbf{Configuration} & \textbf{Result} & \textbf{Interpretation} \\
    \midrule
    Trans (hybrid pass) & Recovers & Different units \\
    Trans (hybrid pass) & Fails    & Same unit \\
    Cis (single pass)   & Works    & Control \\
    \bottomrule
  \end{tabular}}
  \end{center}

  \vspace{5mm}
  Across 16 audited claims, this test has been applied
  \textbf{zero times}. Circuit papers identify component classes
  (``name movers,'' ``S-inhibition heads,'' ``backup name movers'')
  without testing whether those classes are functionally distinct.

  \vspace{4mm}
  \textbf{Cost:} two forward passes + hybrid assembly.

  \vspace{3mm}
  {\footnotesize Genetics. Benzer's cis-trans test (1955), standard
  in functional genomics.}
}
}

\end{columns}

% ════════════════════════════════════════════════════════
% ROW 2 — THREE INTERNAL/EXTERNAL CRITERIA
% ════════════════════════════════════════════════════════
\begin{columns}
\column{0.33}

\block{I6: Double Dissociation}{
  Two interventions that cross: A breaks function X and spares
  function Y; B breaks Y and spares X. One arm alone shows
  involvement. The crossing shows specificity---neither effect
  can be explained by generalized disruption.

  \vspace{4mm}
  \textbf{In MI:} patching a belief-tracking subspace in
  Llama-3-70B yields IIA~$= 1.0$ on belief-framed questions
  and $0.0$ on reality-framed questions about identical
  underlying facts. The subspace responds to the epistemic
  framing, not the information.

  \vspace{4mm}
  Across 16 audited claims, exactly \textbf{one} meets this
  criterion---by different authors (Feucht et~al.\ 2025),
  three years after the original paper.

  \vspace{4mm}
  \textbf{Cost:} one extra condition.

  \vspace{3mm}
  {\footnotesize Neuroscience. Double dissociation logic,
  Shallice (1988).}
}

\column{0.34}

\block{I9: Epistatic Interaction}{
  Do circuit components interact non-additively? Ablate A alone,
  ablate B alone, ablate both. If the joint deficit differs from
  the sum of single deficits, the components interact.

  \vspace{4mm}
  \textbf{In MI:} sub-additive interaction (joint~$<$~sum) means
  redundancy---A and B partially substitute for each other.
  Super-additive interaction (joint~$>$~sum) means synergy---the
  components depend on each other and the circuit is more than
  the sum of its parts.

  \vspace{4mm}
  \begin{center}
  \renewcommand{\arraystretch}{1.3}
  \begin{tabular}{@{}ll@{}}
    \toprule
    \textbf{Pattern} & \textbf{Interpretation} \\
    \midrule
    Joint $=$ sum   & Independent \\
    Joint $<$ sum   & Redundancy \\
    Joint $>$ sum   & Synergy \\
    \bottomrule
  \end{tabular}
  \end{center}

  \vspace{4mm}
  \textbf{Cost:} one extra ablation condition.

  \vspace{3mm}
  {\footnotesize Genetics. Fisher's epistasis; interaction
  structure in gene networks.}
}

\column{0.33}

{
\colorlet{blocktitlebgcolor}{externalpurple}
\block{E5: Graded Response}{
  Partial intervention should produce partial effect. A step
  function means you measured a threshold, not a mechanism.

  \vspace{4mm}
  \textbf{In MI:} sweep ablation strength from 0\% to 100\%
  instead of reporting one on/off number. The shape of the
  curve distinguishes specific mechanisms (graded, monotonic)
  from nonspecific disruption (all-or-nothing collapse at some
  threshold).

  \vspace{4mm}
  Most published circuit evaluations report a single ablation
  point. The curve between zero and complete ablation carries
  more information than either endpoint---and recording it
  costs nothing because the interpolation is already available.

  \vspace{4mm}
  \textbf{Cost:} free.

  \vspace{3mm}
  {\footnotesize Pharmacology. Dose-response curves, Hill
  (1910), receptor theory.}
}
}

\end{columns}

% ════════════════════════════════════════════════════════
% ROW 3 — RESCUE, ONSET/OFFSET, TAKEAWAY
% ════════════════════════════════════════════════════════
\begin{columns}
\column{0.33}

\block{I10: Rescue / Reversibility}{
  Remove a component, observe the deficit, then restore it. If
  function does not recover, the ablation caused collateral
  damage beyond what the mechanism hypothesis predicts.

  \vspace{4mm}
  \textbf{In MI:} re-inject the clean activation after ablation
  and measure recovery. Rescue separates ``this component is
  necessary'' from ``removing this component broke something
  downstream.''

  \vspace{4mm}
  Activation patching already contains the rescue logic: patching
  in a clean activation \emph{is} a rescue experiment. Almost no
  published evaluation reports the recovery rate alongside the
  degradation rate.

  \vspace{4mm}
  \textbf{Cost:} one forward pass.

  \vspace{3mm}
  {\footnotesize Genetics. Rescue experiments in functional
  genomics.}
}

\column{0.34}

\block{I11/I12: Onset and Offset Coupling}{
  A mechanism should appear when its capability appears (onset
  coupling, I11) and disappear when the capability is removed
  (offset coupling, I12). Onset alone is suggestive. Offset
  alone is suggestive. Together they establish a bidirectional
  link between mechanism and capability.

  \vspace{4mm}
  \textbf{In MI:} onset coupling is reported for induction
  heads---the circuit forms at a training phase transition
  when in-context copying emerges. The partner test, offset
  coupling, has not been reported: fine-tune or unlearn the
  capability, then check whether the circuit degrades in
  tandem. If the circuit persists after the capability is
  gone, the causal link runs in one direction only.

  \vspace{4mm}
  \textbf{Cost:} onset = checkpoint tracking (cheap).
  Offset = fine-tuning or unlearning (expensive).

  \vspace{3mm}
  {\footnotesize Medical microbiology. Fredricks \& Relman
  (1996), guidelines for establishing microbial causation.}
}

\column{0.33}

{
\colorlet{blocktitlebgcolor}{colorOne}
\block{Takeaway}{
  {\Large\bfseries Six experiment designs from four source
  disciplines. Most cost one extra condition. Almost none
  have been tried.}

  \vspace{5mm}
  \begin{itemize}[leftmargin=*, itemsep=3mm]
    \item \textbf{1/16} audited claims meets double
      dissociation (I6)
    \item \textbf{0/16} meets complementation (C6)
    \item \textbf{0/16} reports offset coupling (I12)
    \item \textbf{0/16} reports epistatic interaction
      structure (I9)
    \item Dose--response (E5) is free and almost never
      reported
  \end{itemize}

  \vspace{4mm}
  {\small These six criteria are part of a 36-criterion
  framework spanning five validity types. The framework also
  introduces interpretive validity (V1--V5)---the requirement
  that the description mode match the evidence---for which
  no source discipline has a counterpart.}

  \vspace{5mm}
  \begin{center}
  \texttt{elliot@elliottower.ai}

  \vspace{2mm}
  {\small Paper:}
  \texttt{doi.org/10.5281/zenodo.20478480}
  \end{center}
}
}

\end{columns}

\end{document}
